%!TEX root = ../hutbthesis_main.tex

\chapter{研究理论基础}

\section{Model Context Protocol（MCP）协议}

MCP（Model Context
Protocol）是本系统同大模型，仿真器交互的关键通信协议，该协议旨在形成标准化的上下文感知交互机制，达成不同系统之间结构化指令的传递以及状态的同步，FastMCP是MCP协议的一种轻量化达成框架，给本文章赋予了高扩展性的服务端开发根基，它具有如下核心特点以符合人车仿真器交互的需求。

轻量级通信框架：以基于TCP/IP长连接方式进行通讯，实现低延时下发命令以及反馈信息，满足仿真器对实时性需求。

模块化服务设计：以功能划分服务模块（例如车辆控制、编辑器工具等），各个模块可以单独进行注册、启动停止，符合仿真器控制以及引擎编辑的功能拆分理念。

协议有会话上下文保持功能，可以保存仿真器当前情况，如车辆位置、天气等，使得大模型指令根据仿真器当前情况进行生成，避免指令和实际情况不符。

本文章依靠FastMCP框架塑造出HUTB MCP Server和Unreal MCP Server，HUTB MCP
Server关注HUTB/CARLA仿真器的运行时控制，Unreal MCP Server则着眼于Unreal
Engine的编辑器层操作，经由统一的MCP协议接口达成上层交互层（包含语音，Web，AI助手）和底层仿真器之间的解耦通信。

\begin{figure}[H]
\centering
\includegraphics[width=11.8cm,keepaspectratio]{images/image4.png}
\caption{MCP协议完整交互架构图}
\label{fig:mcp-architecture}
\end{figure}

如图~\ref{fig:mcp-architecture}所示，用户通过将自身需求上传到AI客户端，将指令参数发送给MCP服务端，FastMCP工具封装，并将状态结果反馈给AI客户端，用户便可以通过Claude，Web界面获得结论。

\section{人车仿真器底层技术架构}

\subsection{HUTB/CARLA 仿真器核心原理}

HUTB/CARLA仿真器是一个基于Unreal Engine
4开发的开放源代码自动驾驶仿真环境，在本论文中用于车辆、无人机以及环境的物理级仿真，它所依赖的关键技术有：

多智能体交互框架：仿真器内建有可管理多个不同类型的Actor（比如车辆或无人机等）实例以及它们的状态信息，每个Actor都有一个唯一的ID以及自身的物理属性（例如其质量和动力学特性），可以利用接口来创建、销毁或者改变这些Actor的状态。

物理引擎适配：引入PhysX物理引擎，真实表现汽车动力学特性（油门、转向、制动等）及无人机飞行状态（悬停、降落时空气动力学效果）还有天气对于物理世界影响（例如下雨时路面上摩擦力减少）。

传感器仿真体系能够对相机，激光雷达，毫米波雷达等多种传感器实施数字化仿真，传感器采集的数据可以即时反馈到MCP服务端，给大模型给予环境感知的依照，该数据格式同真实硬件输出相契合，保证了控制指令具有现实映射性。

车辆输出状态由油门、转向、制动、自动驾驶状态共同决定，表达式为：

\begin{equation}
Vehicle_{state}=f(Throttle,Steer,Brake,Autopilot)
\label{eq:vehicle-state}
\end{equation}
式中，$Throttle$代表油门控制量，$Steer$表示转向角，$Brake$显示制动强度，$Autopilot$是自动驾驶使能标志。

\begin{figure}[H]
\centering
\includegraphics[width=0.96\linewidth,keepaspectratio]{images/image6.png}
\caption{人车仿真器层次结构原理图}
\label{fig:simulator-hierarchy}
\end{figure}

如图~\ref{fig:simulator-hierarchy}所示，人车仿真器层次结构原理图，该系统采用五层堆栈式架构，自上而下依次为：API控制层，逻辑控制层，物理引擎层，环境与场景层，渲染层。

API控制层是：提供对外接口，使外界程序或者用户可以使用API来控制仿真器，支持多种类型的实体进行控制，可以设置环境参数（如天气等）以及传感器参数，在编辑器中也可以对场景或者参数进行在线调整。

逻辑控制层主要负责：智能体管理仿真中所有自主实体（车辆、行人等）的生命周期以及行为调度。指令解析：将来自上层API高层指令转化为可以被本系统执行的操作。

物理引擎层主要为：求解汽车动力学方程（加速度、速度、位移）。摩擦系数：对汽车轮胎与路面之间以及车身与空气之间的相互作用起着至关重要的作用，从而影响驾驶的真实感。

环境与场景层主要是：提供道路网络、地形、建筑物等静态场景。模拟不同时间、天气下的光照条件。雨、雪、雾等动态气象条件（与API层联动）。其他车辆、行人、骑行者等动态元素

\subsection{Unreal Engine编辑器扩展机制}

Unreal Engine
4的编辑器扩展机制为本研究中的场景编辑以及蓝图编写奠定基础，主要包括：

插件化扩展架构：编写UnrealMCP
Plugin插件，使MCP协议连接到Unreal编辑器核心接口（Editor
Subsystem）上，从而可以对Actor、蓝图编辑、UI控件等进行编程式操作。

蓝图可视化编程体系：Unreal蓝图系统是以节点形式包裹着C++底层代码，而我们利用MCP协议实现蓝图节点的生成、连接以及属性设置等操作，在不需要打开编辑器的情况下就可以进行交互逻辑的设计。

Level Editor主要是Level
Subsystem，这个subsystem提供了对场景中的Actor添加、删除、修改、查询以及移动、旋转、缩放的操作方法，为本文中场景编辑部分提供底层支持。

\section{自然语言交互与仿真器控制的映射逻辑}

自然语言控制主要实现从``自然语言指令到结构化的MCP指令再到仿真器API调用''三者之间的映射关系，而这些都基于以下几点：

指令解析及意图理解：利用大规模预训练模型所具备自然语言处理能力，把未经过任何处理自然语言（如：``让无人机在广场悬停''，``设置暴雨天气''等）转化为结构化意图（操作对象、操作类型及其对应值），比如将``生成一辆轿车在坐标(100,200,0)''转化为``操作对象：车辆，操作类型：生成，参数：蓝图ID=轿车、位置=(100,200,0)、旋转=(0,0,0)''；

API映射规则库：定义自然语言意图到仿真器API间的映射关系，如``自动驾驶''对应set\_vehicle\_autopilot
API，``极端天气设置''对应set\_extreme\_weather
API，规则库包含车辆、无人机、天气、编辑器等方面的所有API间的关系。

参数标准化转换是把自然语言中的模糊参数，例如``广场''、``暴雨''，转化为仿真器能够理解的标准参数，``广场''对应预先定义好的一个位置，``暴雨''则是指有降雨量=100、能见度=50、风速=30等天气参数组成的集合，这样才使得这条命令可以被执行。

\begin{figure}[H]
\centering
\includegraphics[width=0.95\linewidth,keepaspectratio]{images/image7.png}
\caption{自然语言到仿真指令映射流程图}
\label{fig:nl-to-api-flow}
\end{figure}

如图~\ref{fig:nl-to-api-flow}所示，用户将自身需求上传，AI通过对用户需求解析，将其进行参数对比与提取，并将MCP指令化结构，最后通过API执行。
